% Owned by #22 (SOTA metrics table), with results from #20 (MobileNetV2) and
% #21 (EfficientNet-B0). Every number is in tables/sota_comparison.tex,
% generated by scripts/build_tables.py --table sota - none is typed here.

\subsection{Transfer-learning protocol}
\label{sec:sota-protocol}

Two pretrained architectures were adapted to EuroSAT:
MobileNetV2~\cite{sandler2018mobilenetv2}, and EfficientNet-B0~\cite{tan2019efficientnet}.
Both start from ImageNet weights and the final linear layer of each classifier was replaced by a
new ten-class prediction layer. Every run uses the same stratified split as the Models A and B, 
random seed, augmentation, training-split normalisation statistics, batch
size of 128, 30-epoch budget and Colab Tesla~T4 as the machine to run on. 
The checkpoint was chosen by lowest validation loss, before testing.

Once the new prediction layer is in place, there were two cases to choose from, either
keep the pretrained backbone \emph{frozen}, or \emph{fine-tune} it. Each
network was tested using both methods, to decide wether there is a gain by fine tuning all parameters.

\begin{itemize}[leftmargin=*]
  \item \textbf{Frozen.} In this case, only the new prediction layer was trained, with Adam Optimizer at
  a learning rate of $10^{-3}$. The CNN backbone acts as a fixed feature extractor,
  and its batch-normalisation statistics were held at their original ImageNet values.
  \item \textbf{Fine-tuned.} In this case, every weight was trained, starting from the
  ImageNet weights, with Adam at a learning rate of $10^{-4}$. Here we chose a learning rate
  below the frozen-backbone rate to limit changes to the pretrained filters while
  the new classifier learns.
\end{itemize}

Our main runs keep the input at $64\times64$,which is the original resolution that the custom
models see. Because ImageNet networks are
designed for $224\times224$ input, the experiment was repeated for each network
with the $64\times64$ images upsampled to $224\times224$, for both the frozen and fine-tuned cases. Each $224\times224$
pair differs from its original pair in input size only, so as to identify the effect of the original resolution on the pretrained features.

\subsection{Results}
\label{sec:sota-results}

\begin{table}[htbp]
  \centering
  % Nine columns: set small rather than \resizebox-ed, because the Pending
  % box \generatedtable shows before the runs exist cannot sit in a \resizebox.
  \footnotesize
  \setlength{\tabcolsep}{4pt}
  \generatedtable{tables/sota_comparison.tex}{\#22}
  \caption{Pretrained lightweight models on the EuroSAT test split. Precision
  and recall are averages over the ten classes. \emph{Trained} counts the
  parameters updated during training. Sizes are in MB, measured from the saved
  weights; the custom models in Table~\ref{tab:custom-comparison} are in KB.
  The rows under Supplementary section upsample the $64\times64$ images to $224\times224$.}
  \label{tab:sota-comparison}
\end{table}

Table~\ref{tab:sota-comparison} lists the results from both cases for each network. \emph{Parameters} and \emph{Size} are
identical for the frozen and fine-tuned runs of a network, because freezing
changes what is trained, not what is deployed. \emph{Trained} is where they
differ.  In the frozen case it updates only the prediction layer,which is under one
per cent of the network. Accuracy
gaps between the two rows of a pair is therefore bought with training effort
alone, with no cost in memory.

We observed that for both networks that accuracy gap is large. Fine-tuning raises all the metrics by 5--10 points
above the frozen feature extractor CNN case. The explanation is likely the domain shift from ImageNet to Satellite Images.
ImageNet filters were learnt on close-range photographs of objects, while
EuroSAT is overhead satellite images at a much larger scale. 
A frozen backbone can only re-weight features that were learnt for the wrong kind of
image while fine-tuning lets the filters themselves move towards field boundaries,
road textures and water.

The two architectures also rank differently under the two strategies described above. Frozen,
EfficientNet-B0 is clearly the weaker network, despite having almost twice
MobileNetV2's parameters. In the Fine-tuned case, the two networks are practically at the same level. The
gap between them therefore says how well each network's ImageNet features happen to
transfer to $64\times64$ overhead images.
At $64\times64$, both networks' downsampling stages leave a $2\times2$ final
feature map. Different sensitivity of the pretrained features to input
scale is a possible explanation.
Once fine-tuned, however, EfficientNet-B0 offers no meaningful accuracy over
MobileNetV2 for nearly twice the memory. Therefore MobileNetV2 is the better edge
candidate of the two.

The supplementary rows show that upsampling to $224\times224$ gives both frozen backbones a gain of
several points, but the fine-tuned models gain less than one. A possible explanation for this is the fact that frozen
network's filters are fixed at ImageNet's scale, so enlarging the input
is the only way to bring the data to them. Fine-tuning instead moves the
filters to the data, after which resolution matters much less: for each
network, the fine-tuned $64\times64$ model outperforms the frozen
$224\times224$ one. It also does about $(224/64)^2\approx12$ times fewer
operations per image, because a convolution's cost scales with the $H\times W$
of its output. For an edge device,
fine-tuning at native resolution is the cheaper way to get accuracy.

Even the smallest of these models is one to two orders of magnitude larger on
disk than Model~B (Table~\ref{tab:custom-comparison}). Whether the extra accuracy is worth that memory is the
question Section~\ref{sec:task6} answers.

The training curves (Figures~\ref{fig:sota-loss}--\ref{fig:efficientnet-224-loss})
show the split between the two strategies for both networks. Each frozen
run levels off at a clearly higher loss than its fine-tuned partner. At
$64\times64$ the frozen runs' training and validation losses stay close
together. The frozen
EfficientNet-B0 curves are also visibly noisier than MobileNetV2's: after the
first few epochs both of its losses fluctuate instead of
settling, and at $224\times224$ the lowest validation loss happens to fall early,
so its selected checkpoint comes early in the run. The fine-tuned curves of
both networks are smooth and settle within a few epochs.

The noise in the frozen EfficientNet-B0 curves comes from its fixed features,
not from a training fault. While its validation loss swings from epoch to
epoch, its validation accuracy stays almost level, so the predicted classes
barely change; what changes is the confidence given to the many images that
the frozen features do not separate. Cross-entropy penalises a confidently
wrong prediction heavily, so small updates to the prediction layer move the
loss on these borderline images far more than they move the accuracy.
MobileNetV2's frozen features leave fewer such images, and its loss is lower
and flatter. The two backbones also end in different activations.
EfficientNet-B0 uses the unbounded SiLU (Swish)
activation~\cite{tan2019efficientnet}, and a small fraction of images produce
very large pooled features. 
MobileNetV2's ReLU6 caps every feature at
six~\cite{sandler2018mobilenetv2}, the bounded behaviour discussed in
Section~\ref{sec:activation-choice}. In a diagnostic check that retrained only the prediction
layer on fixed features, standardising those features reduced the
fluctuation only slightly. Learning rates of $3\times10^{-4}$ and $10^{-4}$
did not remove it either, so the frozen-head learning rate is not the cause.
The practical cost lies in checkpoint selection. On a noisy validation curve
the lowest-loss epoch is partly a matter of chance, and it need not be the
run's most accurate epoch. Fine-tuning removes the effect because the backbone
adapts its features to EuroSAT, which is one more reason the comparison in
Section~\ref{sec:task6} uses the fine-tuned runs.

The confusion matrices of the fine-tuned models
(Figures~\ref{fig:mobilenetv2-confusion}--\ref{fig:efficientnet-224-confusion})
place the remaining errors in the same few class pairs for both networks.
Highway and River are confused in both directions; at $64\times64$ both are
thin linear features a few pixels wide. The agricultural and vegetation
classes (AnnualCrop, PermanentCrop, Pasture and HerbaceousVegetation), which
differ mainly in texture, account for a large share of the rest.

\begin{figure}[htbp]
  \centering
  \generatedfigure{figures/sota_loss_curves.png}{\#20}
  \caption{Training and validation loss of MobileNetV2 at $64\times64$, with
  the backbone frozen and fine-tuned.}
  \label{fig:sota-loss}
\end{figure}

\begin{figure}[htbp]
  \centering
  \generatedfigure{figures/efficientnet_loss_curves.png}{\#64}
  \caption{Training and validation loss of EfficientNet-B0 at $64\times64$,
  with the backbone frozen and fine-tuned.}
  \label{fig:efficientnet-loss}
\end{figure}

\begin{figure}[htbp]
  \centering
  \generatedfigure{figures/sota_224_loss_curves.png}{\#20}
  \caption{Training and validation loss of MobileNetV2 at $224\times224$ (the
  supplementary resolution ablation), with the backbone frozen and fine-tuned.}
  \label{fig:sota-224-loss}
\end{figure}

\begin{figure}[htbp]
  \centering
  \generatedfigure{figures/efficientnet_224_loss_curves.png}{\#64}
  \caption{Training and validation loss of EfficientNet-B0 at $224\times224$
  (the supplementary resolution ablation), with the backbone frozen and
  fine-tuned.}
  \label{fig:efficientnet-224-loss}
\end{figure}

\begin{figure}[htbp]
  \centering
  \generatedfigure[width=0.82\linewidth]{figures/mobilenetv2_64_ft_confusion_matrix.png}{\#20}
  \caption{Test confusion matrix of the fine-tuned MobileNetV2 at $64\times64$.}
  \label{fig:mobilenetv2-confusion}
\end{figure}

\begin{figure}[htbp]
  \centering
  \generatedfigure[width=0.82\linewidth]{figures/efficientnetb0_64_ft_confusion_matrix.png}{\#64}
  \caption{Test confusion matrix of the fine-tuned EfficientNet-B0 at
  $64\times64$.}
  \label{fig:efficientnet-confusion}
\end{figure}

\begin{figure}[htbp]
  \centering
  \generatedfigure[width=0.82\linewidth]{figures/efficientnetb0_224_ft_b128_confusion_matrix.png}{\#64}
  \caption{Test confusion matrix of the fine-tuned EfficientNet-B0 at
  $224\times224$ (supplementary).}
  \label{fig:efficientnet-224-confusion}
\end{figure}
